\documentclass[final,pdflatex]{sn-jnl}
%\documentclass[submitted,12pt]{article}
\usepackage[font=footnotesize]{caption}
\usepackage{elsnat}
\usepackage{ii}

\newcommand{\TDE}{Transverse Doppler Effect\xspace}
\newcommand{\tDe}{transverse Doppler effect\xspace}
\newcommand{\EP}{\Ehr paradox\xspace}
\newcommand{\mspp}[1]{\mkbibbrackets{p.~#1}}
\newcommand{\igamma}{\ensuremath{\gamma^{-1}}\xspace}

\title{Is Time Dilation \s{Real}? Einstein and the Transverse Doppler Effect}

\author{Marco Giovanelli}\vspace{0.2cm}\affil{\centering
Università di Torino\\
Department of Philosophy and Educational Sciences\\
Via S. Ottavio, 20 10124 - Torino, Italy\\ \vspace{0.2cm}\texttt{marco.giovanelli@unito.it}\vspace{0.5cm}}
\begin{document}



 
% Einstein1949, Einstein1949a: original Schilpp page ranges/title wording.
%Ritz1908b, Pierseaux2003, Bell1976: page ranges need original/PDF confirmation.
%Brown2005, Staley2008, Brown2022: publisher/date/edition discrepancies.
%ECSP: whether to keep combined North-Holland and Interscience publication data.
%EA: no verified publication date; add one only if your style requires an access/date convention.
%Minkowski19078: key likely typo, but renaming requires updating citations.
%Acuña2025: key preserved even though the article is now dated 2026.
 
 
\abstract{
As early as 1907, Einstein realized that the second-order transverse Doppler redshift reported by Stark in the spectral lines of fast-moving ions in canal rays could serve as a direct experimental test of time dilation. He rejected Stark’s \emph{dynamical explanation} of the effect as a real frequency change of the moving ions. Instead, assuming that ions of the same kind emit the same spectral lines when at relative rest, Einstein offered a \emph{kinematical explanation} of the frequency shift as a consequence of time dilation. He labeled the effect \textit{apparent}, since the shifted line is not recorded by a spectrograph co-moving with the ion. At the same time, he regarded it as \textit{real}, in the sense that it would not occur if the old kinematics held, and is thus empirically testable. From the 1920s onward, Einstein argued that the behavior of atomic clocks calls for a dynamical explanation. The present paper shows, however, that even if a special relativistic explanation of the spectral identity of atoms were available, it would merely reinforce the claim that relativistic time dilation admits a purely kinematical interpretation. It concludes that modern dynamicists misread Einstein’s plea for a dynamical account of clock behavior as a demand for an \emph{explanation} of time dilation, whereas it was primarily motivated by the problem of its \emph{confirmation}.}


\maketitle

\section*{Acknowledgements}
Not applicable


\section*{Ethical Statement}
\begin{itemize}
\item  Funding: None \item  Conflict of Interest: None declared
\item  Ethical approval: Not required \item  Informed consent: Not required \item  Author contribution: I am the sole author of this paper
\end{itemize}


\end{document}
